% ============================================================
%  ch07  第二讲 §2.4 Picard 群、§2.5 中心扩张与仿射 Kac--Moody 代数
%  原文行 1652–1765
% ============================================================

\subsection{Picard 群}\label{The Picard group}

为简单起见，设 $G$ 单且单连通。以 $\theta$ 记短支配余权，前面已说明它是 $\xcoch(T)^+-\{0\}$ 的极小元。由命题 \ref{geom of opp Sch}，$\Theta:=\Gr^{\geq \theta}$ 余维为一，且 $\Gr=\Gr^0\sqcup \Theta$。下面说明它是 $\Gr$ 上某个丰沛线丛的整体截面的零点集。

\begin{prop}\label{big cell 2}
存在丰沛线丛 $\mL$ 与整体截面 $\sigma\in\Ga(\Gr,\mL)$，使得 $\Gr^0=\{x\in\Gr \mid \sigma(x)\neq 0\}$。
\end{prop}

\begin{proof}
这里说明想法，细节见 \cite[Lemma 2, Corollary 3]{Fa}。

先考虑 $G=\SL_n$。此时 $L^{<0}G\cdot e$ 是 $\mL_{\det}$ 的某个整体截面 $\sigma$ 的非零集，$\sigma$ 构造如下。令 $L_0=(t^{-1}k[t^{-1}])^n$，考虑 $\Spec R$ 上的两项复形
\[0\to R\otimes L_0\oplus\Lambda\stackrel{\Upsilon}{\to} R\otimes L_0\oplus R[[t]]^n\to 0,\quad \Upsilon(v,w)=(v,0)+\beta(w),\]
其中 $\beta: \La\to R((t))=R\otimes L_0\oplus R[[t]]$ 是自然包含。由定义，该复形（更确切地说，把两项都除以 $ R\otimes L_0\oplus t^NR[[t]]^n$（$N\gg 0$）后所得的复形）的行列式上同调恰是 \S\ref{determinant line} 中构造的行列式线丛 $\mL_{\det}$，而 $\sigma=\det(\Upsilon)$ 定义 $\mL_{\det}$ 的整体截面。把点 $x\in \Gr$ 写成 $g\La_0$（$g\in LG$），则由定义 $\sigma(g\La_0)\neq 0$ 当且仅当 $\Upsilon$ 在该点处是同构。而后一条件又等价于 $g^{-1} L_0=L_0$（见 \cite[Lemma 2]{Fa}）。这就完成了 $\SL_n$ 情形的证明。

其次，取表示 $\rho:G\to \GL(V)$，使得存在 $v\in V$，其在 $\GL(V)$ 中的稳定化子为 $G$（这样的表示存在，因为 $G$ 半单）。设 $f_\rho: \Gr_G\to \GL_{\SL_n}$ 是诱导映射，可以证明 $f_\rho^{-1}(L^{<0}\SL_n\cdot e)=L^{<0}G\cdot e$。于是 $\Theta_G$ 是 $\Theta_{\SL_n}$ 的拉回（作为闭子集）。
\end{proof}

回忆预叠上线丛的定义（见 \S\ref{Intro:sp}）。设 $e\in X(k)$ 是点，定义 $X$ 上在 $e$ 处刚性化的线丛的 Picard 群胚 $\Pic^e(X)$ 为由 $(\mL,\epsilon)$ 构成的群胚，其中 $\mL$ 是 $X$ 上的线丛，$\epsilon: e^*\mL\simeq k$ 是平凡化。注意若 $X$ 是真代数簇，则 $\Pic^e(X)$ 离散，退化为 $X$ 的通常 Picard 群。

\begin{thm}\label{Pic of Gr}
设 $G$ 单且单连通。给 $\Theta$ 赋以既约闭归纳概形结构，则 $\Theta$ 是有效 Cartier 除子，相应的线丛 $\mO(1):=\mO(\Theta)$ 丰沛。此外 $\Pic^e(\Gr)=\bZ\mO(1)$。
\end{thm}

\begin{rmk}\label{generator}
若 $\rho: G\to \GL_n$ 是表示，则 $\mL_{\det}$ 沿 $f_\rho:\Gr_G\to \Gr_{\GL_n}$ 的拉回是 $\mO(1)$ 的正倍数，即 $f_\rho^*\mL_{\det}=\mO(d_\rho)$，其中正整数 $d_\rho$ 称为表示 $\rho$ 的 \textbf{Dynkin 指数}。$d_\rho$ 的一般公式见 \cite[Lemma 5.2]{KNR}，另见 \cite[Lemma 10.6.1]{So} 与 \cite[Lemma 4.2]{YZ}（注意上述文献都只考虑 $\cha\ k=0$，但 \cite[Lemma 4.2]{YZ} 的论证对任意特征都适用）。

例如，若 $\rho$ 是伴随表示，则 $d_\rho=2h^\vee$，其中 $h^\vee$ 是 $G$ 的对偶 Coxeter 数。若 $\rho$ 是 $\SL_n$ 或 $\on{Sp}_{2n}$ 的标准表示，则 $d_\rho=1$。若 $G$ 是正交群 $\on{SO}_n$（设 $\cha\ k>2$）、$\rho$ 是标准表示，或 $G=G_2$、$\rho$ 是唯一的 $7$ 维表示，则 $d_\rho=2$，其唯一的平方根可由 $\on{OGr}(V)$ 上 Pfaffian 线丛 $\mL_{\on{Pf}}$ 的拉回得到（见注 \ref{Pfaffians}）。

换个角度看，上述讨论本质上给出了 $G$ 为典型群或 $G=G_2$ 时 $\mO(1)$ 的构造。但若 $G=E_8$，则对任意 $\rho$ 都有 $60\mid d_\rho$，因此在此情形下找出 $\mO(1)$ 的显式构造是个有趣的问题。
\end{rmk}

\begin{rmk}\label{square root of adjoint}
由此可知，对伴随表示 $\Ad: G\to \GL(\frakg)$，$f_{\Ad}^*\mL^{-1}_{\det}$ 有唯一的刚性化平方根，记作 $\mL_{c}$。这就是仿射 Grassmann 簇上所谓的\textbf{临界线丛}，它在几何朗兰兹纲领中起重要作用（见 \cite{BD}）。
\end{rmk}

证明该定理的一条途径是把 $\Gr$ 等同于相应 Kac--Moody 群的部分旗簇（见下文 \S\ref{Kac-Moody}），然后引用 \cite{Ma} 的结果。下面是更直接的途径：$\mO(1)$ 应当是与除子 $\Theta$ 相配的线丛。但 $\Gr$ 不是有限维的，一般也不能写成光滑有限维子簇的递增并，故其真正的定义依赖下述引理。

\begin{lem}\label{O1 on Bun}
设 $G$ 单且单连通。$\Bun_G(\bP^1)$ 有开子叠 $\mathring{\Bun}_G(\bP^1)$，分类 $\bP^1$ 上可平凡化的 $G$-主丛，且边界 $\Bun_G\setminus \mathring{\Bun}_G$ 纯余维一。因此它是有效 Cartier 除子，在 $\mathring{\Bun}_G(\bP^1)$ 上定义了线丛 $\mO(1)$。
\end{lem}

\begin{proof}
有 $\mathring{\Bun}_G(\bP^1)= [L^-G\backslash L^{<0}G\cdot e]$，而 $\Bun_G\setminus \mathring{\Bun}_G=[L^-G\backslash \Theta]$ 纯余维一。也可以用形变理论证明。
\end{proof}

于是拉回在 $\Gr$ 上定义了线丛，暂记作 $\mL(\Theta)$。注意 $\mL(\Theta)$ 丰沛，因为由命题 \ref{big cell}，它的某个倍数是行列式线丛的拉回。

其次，不难证明 \eqref{curve C} 中构造的有理曲线 $C_{\theta,\theta}$ 在 $t^\theta$ 处与 $\Theta$ 横截相交（例如看切空间），其中 $\theta$ 仍是短支配余根。此外可以证明限制到 $C_{\theta,\theta}$ 诱导 Picard 群的单射（用 Demazure 消解）。于是有包含 $\bZ\mL(\Theta)\subset \Pic^e(\Gr)\subset\Pic(\bP^1)$，它把 $\mL(\Theta)$ 送到 $\mO_{\bP^1}(1)$。定理随之成立。更多细节见 \cite[Corollary 12]{Fa} 与 \cite[\S~10.2]{So}。

下述论断留作练习。

\begin{lem}\label{lifting Aut action}
$\Aut^+(D)$ 在 $\Gr$ 上的作用提升为 $\mO(1)$ 上的唯一作用，它在 $\mO(1)$ 于 $e$ 处的纤维上为平凡作用。
\end{lem}

\subsection{中心扩张与仿射 Kac--Moody 代数}\label{Kac-Moody}

仍设 $k$ 代数闭，$F=k((t))$。回忆仿射 Dynkin 图分类了两类不同的群。设 $\Ga$ 是连通仿射 Dynkin 图，$S$ 是 $\Ga$ 的顶点子集。Bruhat--Tits 的工作对 $\Ga$ 给出 $F$ 上单且单连通的约化群 $G_\Gamma$，对 $S$ 给出 $\mO$ 上的抛孔群概形 $\underline G_S$（概述见 \cite{Ti}，分类见\emph{同上} \S 4.2）。另一方面，对 $\Ga$ 有 Kac--Moody 代数 $\frakg_\Ga$（\cite[Chap. 4]{Kac}）与相应的 Kac--Moody 群 $G_{\Ga}$，对 $S$ 有 $G_{\Ga}$ 的抛物子群 $P_S$\footnote{实际情况要复杂得多：由于 Moody--Teo、Marcuson、Garland、Kac--Peterson、Slodowy、Tits、Mathieu 等人的工作，附着于 $\frakg_\Ga$ 的无穷维群有多种版本（极小的、形式的等等）。概述与文献见 \cite{Ti3}。}。

下面通过解释下图中虚线箭头，说明 Bruhat--Tits 理论与 Kac--Moody 理论的关系。
\begin{center}
\begin{tikzpicture}
  \node (a) at (0,1.7) {$\{\text{$F$ 上的约化群}\}$};
  \node (b) at (-3.2,0) {$\{\text{仿射 Dynkin 图}\}$};
  \node (c) at (0,-1.7) {$\{\text{$k$ 上的仿射 Kac--Moody 代数/群}\}$};
  \draw[dashed, ->] (a) -- (c);
  \draw[->] (b) -- (a);
  \draw[->] (b) -- (c);
\end{tikzpicture}
\end{center}

这里假定 $\Ga$ 是非扭仿射 Dynkin 图\footnote{扭情形更复杂，尤其当 $k$ 特征很小时。}，即 \cite[Chap. 4, TABLE Aff 1]{Kac} 所列的仿射 Dynkin 图。它们可由有限 Dynkin 图（\cite[Chap 4, TABLE Fin]{Kac} 所列）添加一个特殊顶点 $v_0$ 得到。对应于这些 Dynkin 图的 $F$ 上约化群，是 $k$ 上单且单连通约化群作基变换所得。若 $S=\{v_0\}$，则相应的抛孔群概形超特殊，形如 $G\otimes_kk[[t]]$（见例 \ref{hyperspacial and Iwahori}）。

此时，$\Gr_G$ 上的线丛 $\mO(1)$ 可用来定义 $LG$ 被 $\bG_m$ 的中心扩张
\begin{equation}\label{central extension}
1\to\bG_m\to \widehat{LG}\to LG\to 1.
\end{equation}
确切地说，$\widehat{LG}$ 的 $R$-点是 $LG(R)$ 中的 $R$-点 $g$ 连同一个同构
\[g^*\mO(1)\simeq \mO(1).\]
旋转环面的作用典范地提升为 $\widehat{LG}$ 上的作用（由引理 \ref{lifting Aut action}）。令
\[\widetilde{LG}:=\widehat{LG}\rtimes\bG_m^{\on{rot}}.\]

\begin{thm}
归纳群 $\widetilde{LG}$ 同构于由相应非扭仿射 Dynkin 图构造的 Kac--Moody 群 $G_\Ga$ 的（形式版本）。
\end{thm}

这可由下述更精确的定理推出。设 $\underline G$ 是 $G$ 的抛孔群概形，对应于 $G$ 的仿射 Dynkin 图的顶点子集 $S$；设 $P_S\subset G_\Ga$ 是仿射 Kac--Moody 群中相应的抛物子群。则

\begin{thm}
仿射 Grassmann 簇 $\Gr_{\underline G}$ 同构于 \cite{Ma,Ku} 中定义的部分旗簇 $G_\Ga/P_S$。
\end{thm}

\begin{proof}
$G=\SL_n$ 的情形见 \cite{BL}，一般 $G$ 见 \cite{Fa}（另见 \cite[\S 10]{PR}）。想法是先证明 $\Gr_{\underline G}$ 中的 Schubert 簇同构于 Kac--Moody 理论中出现的（先验地不同的）Schubert 簇。
\end{proof}

\begin{rmk}\label{2-cocycle}
$\widehat{LG}$ 的李代数 $\widehat{L\frakg}$ 是 $\frakg(F)$ 被 $k$ 的中心扩张。若 $\cha\ k=0$，它可显式写成 $\frakg\otimes F\oplus k\mathbf{1}$，其李括号由
\[[X\otimes f,\mathbf{1}]=0,\quad [X\otimes f,Y\otimes g]=[X,Y]\otimes fg- (\Res_{t=0}fdg) (X,Y)\mathbf{1}\]
给出，其中 $X,Y\in \frakg$，$f,g\in F$，$(-,-): \frakg\times\frakg\to k$ 是 $\frakg$ 上规范化的不变双线性型，$\Res: Fdt\to k$ 是留数映射。细节见 \cite[\S~7]{Kac}。\zy{关于 nLab：其词条 central extension（\url{https://ncatlab.org/nlab/show/central+extension}，2025-02-16 修订）目前只有导引性内容、未给可引用的定义，故本条不以该词条为依据，而依 \cite[Chap.~4, \S~7]{Kac}：$\widehat{L\frakg}=\frakg(F)\oplus k\mathbf{1}$ 的括号由留数给出，这与 $\widehat{LG}$ 的中心扩张 \eqref{central extension} 在无穷小层面相对应。[来源：A 级 \cite{Kac}；nLab central extension 仅记其现状。]}
\end{rmk}

结合 \cite{Ku,Ma} 的主要结果，可得 Borel--Weil 型定理。为简单起见设 $\cha\ k=0$，并对非扭仿射 Dynkin 图 $\Ga$ 与 $S=\{v_0\}$ 叙述。

\begin{thm}\label{Borel-Weil}
整体截面空间的对偶 $\Ga(\Gr_G,\mO(m))^*$ 同构于 $L_m(\frakg_\Ga)$，即 $\frakg_{\Ga}$ 的第 $m$ 层可积表示。
\end{thm}

Kac--Moody 群的可积表示是约化代数群有限维表示的类比，可用纯表示论方法构造。下面简要回忆 $L_m(\frakg_\Ga)$ 的定义，沿用 \cite{Kac} 的记号。以 $\Lambda_0$ 记 $\frakg_\Ga$ 的基本权（定义见\emph{同上} \S\ 6），以 $M(m\Lambda_0)$ 记 $\frakg_\Ga$ 的最高权 $m\Lambda_0$ 的 Verma 模（定义见\emph{同上} \S\ 9）。则 $L_m(\frakg_\Ga)$ 是 $M(m\Lambda_0)$ 对由\emph{同上} 公式 (10.1.1) 中那些向量生成的子模的商。
